\section*{Giới thiệu}
\addcontentsline{toc}{section}{Giới thiệu}
\textbf{Op-amp} (\textit{Operational Amplifier}) là bộ khuếch đại điện áp có hệ số khuếch đại rất lớn, luôn được dùng cùng điện trở, tụ điện ở mạch hồi tiếp; chính mạch hồi tiếp quyết định chức năng của mạch. Không có hồi tiếp là chế độ \textit{vòng hở} (\textit{open-loop}), chỉ dùng làm bộ so sánh (\textit{comparator}); có hồi tiếp âm là chế độ \textit{vòng kín} (\textit{closed-loop}), làm giảm độ lợi nhưng tăng độ ổn định. Khi hồi tiếp âm, hai ngõ vào đảo ($-$) và không đảo ($+$) có cùng điện thế.

Trong bài này, mô phỏng bằng \textbf{PSpice} các mạch op-amp cơ bản (mục~\ref{sec:follower}--\ref{sec:schmitt}), sau đó thiết kế hai mạch trên \textbf{Altium Designer} (mục~\ref{sec:led}--\ref{sec:relay}).

\subsection*{Cấu hình không đảo}
\given{fig5_01}{0.55\linewidth}{Non inverting configuration\label{fig:5-01}}

\begin{notebox}{Kiểm tra: giải thích $V_{OUT} = 4\,\text{V}$ (Hình~\ref{fig:5-01})}
\dotlines{4}
\end{notebox}

\subsection*{Cấu hình đảo}
\given{fig5_02}{0.55\linewidth}{Inverting configuration\label{fig:5-02}}

\begin{notebox}{Kiểm tra: tính $V_{OUT}$ (Hình~\ref{fig:5-02}), đáp số $-2\,\text{V}$}
\dotlines{4}
\end{notebox}

\begin{notebox}{Quy ước chung khi mô phỏng PSpice}
\begin{itemize}[nosep]
  \item Linh kiện \texttt{OPAMP} lấy từ thư viện \textit{Favorites}; nguồn DC: \texttt{VDC}, nguồn AC/sin: \texttt{VAC} / \texttt{VSIN}; điện trở: \texttt{R}, tụ: \texttt{C}.
  \item \textbf{Bias Point}: đặt ``Enable Bias Voltage / Current Display'' để hiện $V$, $I$ ngay trên sơ đồ.
  \item Lưu ảnh sơ đồ và đồ thị vào \texttt{images/sim/} theo đúng tên tệp ghi ở khung màu cam.
\end{itemize}
\end{notebox}
